\documentclass[12pt,a4paper]{amsart}

\usepackage{amsmath,amssymb,amsthm}
\usepackage{mathtools}
\usepackage[colorlinks=true,linkcolor=blue!60!black,citecolor=green!50!black,
            urlcolor=blue!70!black]{hyperref}
\usepackage[shortlabels]{enumitem}
\usepackage{booktabs}

%% ---- Theorem environments ----------------------------------------
\newtheorem{theorem}{Theorem}[section]
\newtheorem{lemma}[theorem]{Lemma}
\newtheorem{proposition}[theorem]{Proposition}
\newtheorem{corollary}[theorem]{Corollary}
\theoremstyle{definition}
\newtheorem{definition}[theorem]{Definition}
\newtheorem{remark}[theorem]{Remark}
\newtheorem{assumption}[theorem]{Assumption}
\newtheorem{conjecture}[theorem]{Conjecture}
\newtheorem{problem}{Problem}[section]

%% ---- Notation (aligned with Papers I and II) --------------------
\newcommand{\PW}{\mathrm{PW}_\omega}
\newcommand{\Kop}{\mathcal{K}_c}
\newcommand{\GVM}{G^{(N)}_{\mathbf{p},c}}
\newcommand{\psin}{\psi_n^{(c)}}
\newcommand{\EN}{E_{N,c}}
\newcommand{\PN}{P_{N,c}}
\newcommand{\norm}[1]{\|#1\|}
\newcommand{\ip}[2]{\langle #1,\, #2\rangle}
\newcommand{\FN}{\mathcal{F}_N}
\newcommand{\Eout}{\mathcal{E}_{\mathrm{out}}}
%% [PATCH C] Replaced \muhat by \mumN for notational clarity
\newcommand{\mumN}{\mu_{mn}}
\newcommand{\Emn}{\mathbb{E}_{mn}}
\newcommand{\rhom}{\rho_m}
\newcommand{\rhon}{\rho_n}
\newcommand{\rhocl}{\rho_n^{\mathrm{cl}}}

%% ---- Paper metadata ---------------------------------------------
\title[PSWF Product Spectral Tail Estimates]
      {PSWF Product Spectral Tail Estimates:\\
       Bulk and Off-Diagonal Regimes,\\
       and an Obstruction at the Spectral Edge\\[4pt]
       \large (Paper~III in the Prolate--Weil Program)}

\thanks{This paper is the third in the Prolate--Weil program.
Paper~I \cite{PaperI} established frame coercivity of the prolate Gram
form under the Discrete Spectral Transfer Property (DSTP) of \cite{PaperI}.
The companion paper \cite{PaperII_quad} organised the verification of
the DSTP for XRY PSWF-Gauss quadrature into two explicit
conjectural inputs: a PSWF product spectral tail conjecture and an
XRY stability conjecture. The present paper addresses the first of
these two inputs.
In the bulk regime, the tail bound is proved conditional on
Assumption~\ref{ass:bulkconv} (exponential out-of-band convolution
decay); this is identified as an open problem (Problem~\ref{prob:bulkconv}).
A conditional bulk decorrelation lemma (Lemma~\ref{lem:bulk-reduction})
reduces the energy equidistribution problem to weak convergence of
$\psi_n^2$ (Problem~\ref{prob:pswf-weak-limit}).
In the off-diagonal regime, an unconditional uniform bound is proved
(Theorem~\ref{thm:offdiag}), mean spectral localization is established
unconditionally (Proposition~\ref{prop:mean-loc}), and an exact energy
decomposition with an unconditional lower bound is proved
(Proposition~\ref{prop:spectral-lower}); the stronger pointwise
spectral localization is left as an open problem (Problem~\ref{prob:cg}).
The edge regime $m,n \sim N$ is identified as a structural obstruction.}

\author{Anonymous Author}
\address{}
\email{}
\date{\today}

\subjclass[2020]{42C05, 42B10, 47B32, 41A60, 47G10}

\keywords{prolate spheroidal wave functions, spectral tail estimates,
product family, Slepian energy identity, out-of-band energy,
band-edge convolution, frame stability, quadrature compatibility,
bulk convolution decay, spectral localization, mean spectral localization,
energy decomposition, semiclassical analysis,
bulk decorrelation, weak convergence, classical density,
PSWF Clebsch--Gordan, edge regime obstruction, bandwidth doubling}

\begin{document}
\maketitle

\begin{abstract}
Let $\{\psin\}_{n \ge 0}$ be the prolate spheroidal wave functions
(PSWFs) on $[-T,T]$ with bandwidth parameter $\omega > 0$ and
time-bandwidth product $c = \omega T$. Let $\PN$ denote the
orthogonal projector in $L^2([-T,T])$ onto the prolate subspace
$V_{N,c} = \mathrm{span}\{\psi_0^{(c)}, \ldots, \psi_{N-1}^{(c)}\}$.
The companion paper \cite{PaperII_quad} requires a uniform tail bound
for the product family
$\FN := \{\psi_m^{(c)}\overline{\psin} : 0 \le m,n \le N\}$.

The central analytical tool is the \emph{Slepian energy identity}
(Lemma~\ref{lem:slepian}), which bounds the spectral tail
$\|(I-\PN)f\|_{L^2([-T,T])}^2$ by the out-of-band Fourier
energy $\Eout(f) := \int_{|\xi|>\omega}|\hat{f}(\xi)|^2\,d\xi$.
For $f_{mn} = \psi_m^{(c)}\overline{\psin}$, the convolution
$\widehat{f_{mn}} = \hat\psi_m^{(c)} * \hat\psi_n^{(c)}$
is supported in $[-2\omega,2\omega]$ (bandwidth doubling).

We prove:
\begin{enumerate}[(I)]
  \item \textbf{Bulk regime} ($m, n \le \gamma N$, $\gamma < 1$):
  $\|(I-\PN)f_{mn}\|_{L^2} \le C_\gamma e^{-\alpha_\gamma N}$,
  conditional on Assumption~\ref{ass:bulkconv}
  (Problem~\ref{prob:bulkconv}).
  A bulk decorrelation reduction is proved unconditionally
  (Lemma~\ref{lem:bulk-reduction}); the hypothesis of that lemma
  is supplied unconditionally by Paper~IV \cite{PaperIV},
  Theorem~1 (weak convergence of PSWF densities to the classical
  equilibrium measure, with rate $O(1/n)$).
  \item \textbf{Off-diagonal regime} ($|m-n| \ge \delta N$, $\delta > 0$):
  \begin{enumerate}[(a)]
    \item \emph{Unconditional uniform bound:}
    $\|(I-\PN)f_{mn}\|_{L^2} \le C_\delta T^{1/2}$
    (Theorem~\ref{thm:offdiag}).
    \item \emph{Unconditional mean spectral localization:}
    $\Emn[\chi_k] = \mumN + E_{mn}$
    (Proposition~\ref{prop:mean-loc}).
    \item \emph{Exact energy decomposition and lower bound:}
    $\Emn[\chi_k] \ge \mumN/2$, with an explicit energy formula
    (Proposition~\ref{prop:spectral-lower}).
    \item \emph{Conditional algebraic decay:}
    $\|(I-\PN)f_{mn}\|_{L^2} \le C_p(1+|m-n|)^{-p}$
    under Assumption~\ref{ass:specloc}
    (Problem~\ref{prob:cg}).
  \end{enumerate}
  \item \textbf{Edge regime obstruction} ($m, n \sim N$):
  $\Eout(f_{nn}) \ge c_0 > 0$
  (Proposition~\ref{prop:bwdoubling}).
\end{enumerate}
No claim is made about zeros of $\zeta(s)$.
\end{abstract}

\tableofcontents

%%=================================================================
\section{Introduction}\label{sec:intro}
%%=================================================================

\subsection{Context and motivation}

The Prolate--Weil program \cite{PaperI,PaperII,PaperII_quad} seeks a
spectral realization of the Weil positivity criterion
\cite{Weil1952,CCM2025} via the prolate concentration operator
$\Kop$ and its associated Gram forms.
In \cite{PaperI}, the prolate Gram form was shown to be coercive
under the Discrete Spectral Transfer Property (DSTP); see
Section~\ref{ssec:context}.
The companion paper \cite{PaperII_quad} reduced verification of
the DSTP to two conjectural inputs; see
Section~\ref{ssec:context}.

The first input is a uniform tail bound for the product family
$f_{mn} = \psi_m^{(c)}\overline{\psin}$.
This paper establishes such bounds in the bulk and off-diagonal
regimes, proves mean spectral localization and an exact energy
decomposition unconditionally, identifies the structural obstruction
at the edge, and isolates explicit open problems.

\subsection{Strategy and new results}

The Slepian energy identity (Lemma~\ref{lem:slepian})
reduces to bounding $\Eout(f_{mn})$ in the bulk.
In the off-diagonal regime, the near-eigenvector structure
(Section~\ref{ssec:neareig})
\begin{equation}\label{eq:near-eig-intro}
  D_c f_{mn} = \mumN f_{mn} + R_{mn},
  \qquad \mumN := \chi_m(c) + \chi_n(c),
\end{equation}
yields the \emph{mean spectral localization} identity
(Proposition~\ref{prop:mean-loc}).
A subsequent integration-by-parts analysis
(Lemma~\ref{lem:ibp-exact}, Corollary~\ref{cor:emn-explicit})
replaces the commutator error $E_{mn}$ by an
\emph{explicit, fully positive energy functional}:
\begin{equation}\label{eq:energy-intro}
  \Emn[\chi_k]
  = \frac{\mumN}{2}
  + \frac{\frac{1}{2}\int_{-T}^T p(x)
    \bigl[(\psi_m')^2\psi_n^2 + (\psi_n')^2\psi_m^2\bigr]\,dx
    + \frac{c^2}{T^2}\int_{-T}^T x^2\psi_m^2\psi_n^2\,dx}
  {\|f_{mn}\|_{L^2}^2},
\end{equation}
where $p(x) = 1 - x^2/T^2$.
Since every term on the right is non-negative,
$\Emn[\chi_k] \ge \mumN/2$ unconditionally
(Proposition~\ref{prop:spectral-lower}).

New in this version: Lemma~\ref{lem:bulk-reduction} shows that
the energy equidistribution goal $A_{mn}+B_{mn}\approx\mumN/2$
reduces, via a covariance decomposition, to the weak convergence
of $\psi_n^2$ to its classical limiting density; this is
Problem~\ref{prob:pswf-weak-limit}.
That problem is now resolved unconditionally by Paper~IV
\cite{PaperIV} (Theorem~1, with explicit rate $O(1/n)$),
closing the bulk decorrelation step without any WKB assumption.

\subsection{Main results summary}

\begin{theorem}[Bulk; Theorem~\ref{thm:bulk}]
Under Assumption~\ref{ass:bulkconv}, for $m,n \le \gamma N$:
$\|(I-\PN)f_{mn}\|_{L^2} \le C_\gamma e^{-\alpha_\gamma N}$.
\end{theorem}

\begin{theorem}[Off-diagonal, unconditional; Theorem~\ref{thm:offdiag}]
For $|m-n| \ge \delta N$:
$\|(I-\PN)f_{mn}\|_{L^2} \le C_\delta T^{1/2}$.
\end{theorem}

\begin{proposition}[Mean localization; Proposition~\ref{prop:mean-loc}]
For all $m,n \le N$:
$\Emn[\chi_k] = \mumN + E_{mn}$,
where $E_{mn}$ is given by~\eqref{eq:energy-intro} minus $\mumN$.
\end{proposition}

\begin{proposition}[Spectral lower bound; Proposition~\ref{prop:spectral-lower}]
For all $m,n \le N$ with $f_{mn} \not\equiv 0$:
$\Emn[\chi_k] \ge \mumN/2$.
\end{proposition}

\begin{theorem}[Off-diagonal, conditional; Theorem~\ref{thm:offdiag-strong}]
Under Assumption~\ref{ass:specloc}, for $|m-n| \ge \delta N$:
$\|(I-\PN)f_{mn}\|_{L^2} \le C_p(1+|m-n|)^{-p}$.
\end{theorem}

\begin{lemma}[Bulk decorrelation reduction; Lemma~\ref{lem:bulk-reduction}]
For $m,n \le \gamma N$, the energy equidistribution goal
$A_{mn}+B_{mn} = \mumN/2 + o(\mumN)$
follows from weak convergence of $\psi_n^2$ to
$\lambda_n(c)\,\rhocl$; see Problem~\ref{prob:pswf-weak-limit}.
This hypothesis is supplied unconditionally by Paper~IV \cite{PaperIV}.
\end{lemma}

\subsection{Relation to companion paper}\label{ssec:context}

In \cite{PaperII_quad}, verification of the DSTP of \cite{PaperI}
was reduced to:
\begin{enumerate}[(C1)]
  \item PSWF product spectral tail (Conjecture~3.1 of \cite{PaperII_quad}).
  \item XRY stability (Conjecture~4.1 of \cite{PaperII_quad}).
\end{enumerate}
This paper addresses (C1) unconditionally in the mean and energy
localization sense, and conditionally in the tail sense.
(C2) is not addressed.

\subsection{What this paper does not claim}

No statement is made about zeros of $\zeta(s)$.
Assumptions~\ref{ass:bulkconv} and~\ref{ass:specloc} are not proved;
both are formulated as explicit open problems.

\subsection{Organization}

Section~\ref{sec:prelim} establishes PSWF conventions and assumptions.
Section~\ref{sec:tools} gives the Slepian identity and convolution estimate.
Section~\ref{sec:bulk} proves the bulk bound and the decorrelation reduction.
Section~\ref{sec:offdiag} proves off-diagonal bounds, mean localization,
and the exact energy decomposition.
Section~\ref{sec:edge} analyzes the edge obstruction.
Section~\ref{sec:open} collects open problems.

%%=================================================================
\section{Preliminaries}\label{sec:prelim}
%%=================================================================

\subsection{PSWF conventions}

Fix $\omega > 0$, $T > 0$, $c = \omega T > 0$.
Fourier convention: $\hat{f}(\xi) = \int_{\mathbb{R}} f(x)e^{-2\pi i\xi x}\,dx$.

\begin{remark}[Operator conventions]\label{rem:operator-conventions}
Throughout this paper, $c$ is fixed and we write $\psi_n^{(c)}$
in formal statements to emphasise dependence on the bandwidth
parameter. In proofs where $c$ is fixed, we abbreviate
$\psi_n^{(c)} = \psi_n$.
The Sturm--Liouville operator $D_c$ (defined in
Section~\ref{ssec:SL} below) plays the role of a
\emph{quantum Hamiltonian} in the semiclassical analysis of
Sections~\ref{sec:bulk} and~\ref{sec:offdiag}: its eigenvalues
$\chi_n(c)$ are the energy levels, and the spectral mean
$\Emn[\chi_k]$ is the quantum-mechanical expectation value of
$D_c$ in the state $f_{mn}$.
\end{remark}

The \emph{prolate concentration operator}:
\begin{equation}\label{eq:Kop_def}
  \Kop = R_T \circ B_\omega \circ R_T,
\end{equation}
where $B_\omega = \mathcal{F}^{-1}\mathbf{1}_{[-\omega,\omega]}\mathcal{F}$
and $(R_T f)(x) = f(x)\mathbf{1}_{[-T,T]}(x)$.
Eigenvalues $1 > \lambda_0(c) > \lambda_1(c) > \cdots > 0$,
eigenfunctions $\{\psin\}_{n\ge 0}$ orthonormal in $L^2(\mathbb{R})$,
$\ip{\psi_m^{(c)}}{\psin}_{L^2([-T,T])} = \lambda_n(c)\delta_{mn}$.

\subsection{Time concentration and leakage}

\begin{equation}\label{eq:time_conc}
  \lambda_n(c) = \|\psin\|_{L^2([-T,T])}^2,
  \qquad
  \ell_n := 1-\lambda_n(c)
  = \|\psin\|_{L^2(\mathbb{R}\setminus[-T,T])}^2.
\end{equation}

\begin{remark}[Leakage is a time-domain quantity]\label{rem:leakage}
$\ell_n$ measures $L^2$-mass of $\psin$ outside $[-T,T]$ only.
$\int_{|\xi|>\omega}|\hat\psin|^2 = 0$ exactly
(Section~\ref{ssec:bandlimited}).
\end{remark}

\subsection{Bandlimited structure}\label{ssec:bandlimited}

\begin{equation}\label{eq:bandlimited}
  \operatorname{supp}\hat\psin \subset [-\omega,\omega],
  \qquad \|\hat\psin\|_{L^2}^2 = 1.
\end{equation}

\subsection{Product family and bandwidth doubling}

\begin{definition}\label{def:product_family}
$f_{mn} := \psi_m^{(c)}\overline{\psin}$;
$\FN := \{f_{mn}: 0 \le m,n \le N\}$.
\end{definition}

\begin{equation}\label{eq:product_fourier}
  \widehat{f_{mn}} = \hat\psi_m^{(c)} * \hat\psin,
  \quad \operatorname{supp}\widehat{f_{mn}} \subset [-2\omega,2\omega],
  \quad f_{mn} \in \mathrm{PW}_{2\omega}.
\end{equation}

\subsection{Sturm--Liouville operator and spectral identities}\label{ssec:SL}

Each $\psin$ satisfies:
\begin{equation}\label{eq:SL}
  D_c\psin = \chi_n(c)\psin,
  \qquad D_c = -\frac{d}{dx}\!\left(p(x)\frac{d}{dx}\right) + \frac{c^2 x^2}{T^2},
  \quad p(x) := 1 - \frac{x^2}{T^2},
\end{equation}
where $D_c$ is self-adjoint on $L^2([-T,T])$ with
\begin{equation}\label{eq:chi-asym}
  \chi_n(c) = n(n+1) + \tfrac{1}{2}c^2 + O(n^{-1}c^2),
  \qquad n\to\infty.
\end{equation}
See \cite[Ch.~2]{OsipovRokhlinXiao}.
Multiplying~\eqref{eq:SL} by $\psin$ and integrating by parts gives
the \emph{SL energy identity}:
\begin{equation}\label{eq:SL-energy}
  \int_{-T}^T p(x)\bigl(\psi_n^{(c)}\bigr)'^2\,dx
  + \frac{c^2}{T^2}\int_{-T}^T x^2\bigl(\psin\bigr)^2\,dx
  = \chi_n(c)\,\lambda_n(c),
\end{equation}
using $p(\pm T) = 0$ so boundary terms vanish.

\begin{remark}[Form domain of $D_c$ and regularity of $f_{mn}$]
\label{rem:form-domain}
The quadratic form associated to $D_c$ is
\[
  q_{D_c}(u,u)
  = \int_{-T}^T p(x)|u'(x)|^2\,dx
  + \frac{c^2}{T^2}\int_{-T}^T x^2|u(x)|^2\,dx,
\]
with form domain $\mathcal{D}(q_{D_c})
= \{u \in H^1([-T,T]) : \sqrt{p}\,u' \in L^2\}$.
Since $\psi_m^{(c)}, \psin \in C^\infty([-T,T])$
(\cite[Ch.~2, Thm.~2.1]{OsipovRokhlinXiao}),
the product $f_{mn} = \psi_m^{(c)}\psin$ lies in
$C^\infty([-T,T]) \subset \mathcal{D}(q_{D_c})$.
In particular, all integration-by-parts steps in
Lemma~\ref{lem:ibp-exact} are justified without approximation
arguments: boundary terms at $x = \pm T$ vanish exactly because
$p(\pm T) = 0$, independently of the values of $f_{mn}$ at the
endpoints.
\end{remark}

\subsection{PSWF coefficient sequence}

\begin{definition}\label{def:coeff}
$a_k^{mn} := \ip{f_{mn}}{\psi_k^{(c)}}_{L^2([-T,T])}$.
Spectral tail:
$\|(I-\PN)f_{mn}\|_{L^2}^2 = \sum_{k \ge N}|a_k^{mn}|^2$.
\end{definition}

\begin{definition}[Spectral mean]\label{def:spectral-mean}
\begin{equation}\label{eq:spectral-mean-def}
  \Emn[\chi_k]
  := \frac{\displaystyle\sum_{k=0}^{\infty}\chi_k(c)\,|a_k^{mn}|^2}
           {\|f_{mn}\|_{L^2([-T,T])}^2}
  = \frac{\ip{D_c f_{mn}}{f_{mn}}_{L^2}}{\|f_{mn}\|_{L^2}^2}.
\end{equation}
This is the exact spectral expectation value of $D_c$ in state $f_{mn}$.
\end{definition}

\subsection{Spectral projector and index regimes}

$\PN f := \sum_{k=0}^{N-1}\ip{f}{\psi_k^{(c)}}_{L^2}\psi_k^{(c)}$.

\begin{center}
\renewcommand{\arraystretch}{1.3}
\begin{tabular}{lll}
\toprule
Regime       & Condition                            & Section \\
\midrule
Bulk         & $m,n \le \gamma N$, $\gamma\in(0,1)$ & \S\ref{sec:bulk}    \\
Off-diagonal & $|m-n| \ge \delta N$, $\delta>0$     & \S\ref{sec:offdiag} \\
Edge         & $m,n \ge (1-\delta)N$                & \S\ref{sec:edge}    \\
\bottomrule
\end{tabular}
\end{center}

\subsection{Assumptions}

\begin{assumption}[Landau--Widom]\label{ass:lw}
For $N \le (2/\pi-\varepsilon)c$, there exist
$C(\varepsilon),\eta(\varepsilon) > 0$ such that
$\ell_n \le Ce^{-\eta c}$ for all $n \le N-1$.
In particular $\lambda_{N-1}(c) > 1/2$.
\end{assumption}

See \cite[Thm.~1]{Widom1964} and \cite[Ch.~4]{OsipovRokhlinXiao}.

\begin{assumption}[Bulk convolution decay]\label{ass:bulkconv}
For fixed $\gamma\in(0,1)$, there exist
$C_{\mathrm{bc}}(\gamma),\alpha_{\mathrm{bc}}(\gamma)>0$ such that
for all $m,n\le\gamma N$ and $c\ge c_0(\gamma)$:
\[
  \Eout(f_{mn}) \le C_{\mathrm{bc}}\,e^{-\alpha_{\mathrm{bc}}c}.
\]
\end{assumption}

\begin{remark}[Assumption~\ref{ass:bulkconv} is not circular]
\label{rem:bulkconv-not-circular}
Assumption~\ref{ass:bulkconv} is a Fourier-side statement
($\Eout$ measures out-of-band Fourier energy of $f_{mn}$),
whereas the bulk tail bound of Theorem~\ref{thm:bulk} is an
$L^2$-projection statement ($\|(I-\PN)f_{mn}\|_{L^2}$).
The Slepian identity (Lemma~\ref{lem:slepian}) translates between
the two sides; Assumption~\ref{ass:bulkconv} is therefore
strictly weaker than the conclusion.
Evidence: \cite[Ch.~6]{OsipovRokhlinXiao} and
\cite[Eq.~(4.3)]{Slepian1978} give pointwise Fourier decay for
individual PSWFs near the band edge; the convolution stability
for products is the remaining gap (Problem~\ref{prob:bulkconv}).
\end{remark}

\begin{assumption}[PSWF spectral localization]\label{ass:specloc}
For fixed $\delta\in(0,1)$ and $p\ge 1$, there exist $C_p(\delta)>0$
such that for all $|m-n|\ge\delta N$, $0\le m,n\le N$, and $k\ge N$:
\begin{equation}\label{eq:specloc}
  |a_k^{mn}|
  \le \frac{C_p(\delta)}{(1+|\chi_k(c)-\mumN|)^p},
  \qquad \mumN := \chi_m(c)+\chi_n(c).
\end{equation}
\end{assumption}

\begin{remark}[Geometric motivation for Assumption~\ref{ass:specloc}]
\label{rem:specloc-geometric}
Assumption~\ref{ass:specloc} is not merely a reformulation of the
desired conclusion: it has independent geometric support in the
off-diagonal regime.
Lemma~\ref{lem:gap-od} shows unconditionally that for $|m-n|\ge\delta N$
and $k\ge N$, the spectral gap satisfies
$|\chi_k(c)-\mumN|\ge c_3(\delta)|m-n|^2$;
this gap grows with $|m-n|$ and is therefore larger in precisely
the regime where Assumption~\ref{ass:specloc} is claimed.
Moreover, Proposition~\ref{prop:mean-loc} establishes that the
\emph{first moment} of $|a_k^{mn}|^2$ is concentrated near
$\chi_k = \mumN$, consistent with the pointwise decay asserted by
Assumption~\ref{ass:specloc}.
The gap between the first-moment information and the pointwise
localization is the content of Problem~\ref{prob:cg}:
Assumption~\ref{ass:specloc} is not vacuous, but it requires
controlling commutator iterates, which is where the difficulty lies.
\end{remark}

\begin{remark}[Bulk--Off-diagonal connection]
\label{rem:bulk-od-connection}
For $m,n \le \gamma N$ with $\gamma < 1/\sqrt{2}$,
$k^* < N$; hence Assumption~\ref{ass:specloc} implies
exponential tail decay.
\end{remark}

%%=================================================================
\section{Analytical Tools}\label{sec:tools}
%%=================================================================

\subsection{Out-of-band energy and Slepian identity}

\begin{definition}\label{def:eout}
$\Eout(f) := \int_{|\xi|>\omega}|\hat f(\xi)|^2\,d\xi$.
\end{definition}

\begin{lemma}[Slepian energy identity]\label{lem:slepian}
For every $f\in L^2(\mathbb{R})$,
with $a_k = \ip{f}{\psi_k^{(c)}}_{L^2([-T,T])}$:
\begin{align}
  \textstyle\sum_k\lambda_k(c)|a_k|^2
  &= \textstyle\int_{-\omega}^{\omega}|\hat f|^2\,d\xi,
  \label{eq:sei-in}\\
  \textstyle\sum_k(1-\lambda_k(c))|a_k|^2
  &= \Eout(f),\label{eq:sei-out}
\end{align}
and consequently
$\|(I-\PN)f\|_{L^2}^2 \le \Eout(f)/(1-\lambda_{N-1}(c))$.
\end{lemma}

\subsection{Band-edge convolution estimate}

\begin{lemma}[Support geometry]\label{lem:geom}
For $g_m,g_n$ supported in $[-\omega,\omega]$ and
$\xi\in(\omega,\omega+\delta)$:
$(g_m*g_n)(\xi) = \int_{\xi-\omega}^{\omega}g_m(\xi-\eta)g_n(\eta)\,d\eta$;
every contributing $\eta$ satisfies $|\eta|\in(\omega-\delta,\omega)$.
\end{lemma}

\begin{lemma}[Band-edge bound]\label{lem:bandedge}
For PSWFs: $\Eout(f_{mn}) \le 4\omega$.
\end{lemma}

%%=================================================================
\section{Bulk Regime}\label{sec:bulk}
%%=================================================================

\begin{theorem}[Bulk spectral tail bound]\label{thm:bulk}
Under Assumptions~\ref{ass:lw} and~\ref{ass:bulkconv},
for $m,n \le \gamma N$ and $c$ sufficiently large:
$\|(I-\PN)f_{mn}\|_{L^2} \le C_\gamma e^{-\alpha_\gamma N}$.
\end{theorem}

\begin{proof}
Slepian: $\|(I-\PN)f_{mn}\|^2 \le \Eout(f_{mn})/(1-\lambda_{N-1})
\le 2C_{\mathrm{bc}}e^{-\alpha_{\mathrm{bc}}c}
\le 2C_{\mathrm{bc}}e^{-\alpha_\gamma N}$.
\end{proof}

\subsection{Bulk decorrelation reduction}\label{ssec:bulk-reduction}

We now connect the energy equidistribution goal of
Problem~\ref{prob:comm-refined} to a single classical problem:
weak convergence of $\psi_n^2$ to the classical equilibrium density.

\begin{definition}[Normalised PSWF densities]\label{def:densities}
For $0 \le n \le N$, define the probability densities on $[-T,T]$:
\[
  \rhon(x) := \frac{\psi_n^{(c)}(x)^2}{\lambda_n(c)}, \qquad
  \rhom(x) := \frac{\psi_m^{(c)}(x)^2}{\lambda_m(c)}.
\]
The \emph{product density} is
\[
  \rho_{mn}(x)
  := \frac{f_{mn}(x)^2}{\|f_{mn}\|_{L^2([-T,T])}^2}
  = \frac{\psi_m(x)^2\psi_n(x)^2}{\|\psi_m\psi_n\|_{L^2}^2}.
\]
\end{definition}

\begin{definition}[Classical equilibrium density]\label{def:classical}
For $n \le \gamma N$ with $\gamma < 1$, the Sturm--Liouville
potential of $D_c$ is $V(x) = c^2x^2/T^2$ and the eigenvalue is
$E_n := \chi_n(c)$.
The \emph{classical turning points} are
$x_{\pm}(n) = \pm T\sqrt{E_n}/c$,
and the \emph{classical equilibrium density} is
\begin{equation}\label{eq:classical-density}
  \rhocl(x)
  := \frac{1}{\pi}\frac{1}{\sqrt{E_n/c^2 - x^2/T^2}\cdot T}
  = \frac{c}{\pi\,T\sqrt{E_n - c^2x^2/T^2}}
  \qquad (|x| < x_+(n)),
\end{equation}
normalised so that $\int_{-x_+(n)}^{x_+(n)}\rhocl(x)\,dx = 1$.
\end{definition}

\begin{remark}[WKB heuristic]\label{rem:wkb-heuristic}
For large $n$ and $x$ in the interior of $(-x_+(n),x_+(n))$,
the WKB approximation gives
$\psin(x) \approx A_n(x)\cos(\phi_n(x))$
with amplitude $A_n(x) \asymp (E_n - V(x))^{-1/4}$.
Squaring and averaging over the rapid oscillations of
$\cos^2(\phi_n(x)) \approx 1/2$:
$\psi_n(x)^2 \approx \tfrac{1}{2}A_n(x)^2 \asymp (E_n-V(x))^{-1/2}$.
After normalisation, $\psi_n(x)^2\,dx \rightharpoonup \lambda_n(c)\,\rhocl(x)\,dx$.
This heuristic is made rigorous by Paper~IV \cite{PaperIV},
Theorem~1, via Pr\"{u}fer phase monotonicity and a
1D Riemann--Lebesgue argument, with explicit rate $O(\lambda_n/n)$
for $f \in C^1([-T,T])$.
\end{remark}

\begin{lemma}[Bulk decorrelation reduction]\label{lem:bulk-reduction}
Let $m,n \le \gamma N$ with $\gamma < 1/2$.
Suppose that there exists a modulus $\omega_n(\cdot)$ such that
for all $f \in C^1([-T,T])$:
\begin{equation}\label{eq:weak-conv-hyp}
  \left|
    \int_{-T}^T f(x)\,\psi_n(x)^2\,dx
    - \lambda_n(c)\int_{-T}^T f(x)\,\rhocl(x)\,dx
  \right|
  \le \lambda_n(c)\,\omega_n(\|f'\|_{L^\infty}).
\end{equation}
Then the potential energy term in~\eqref{eq:energy-intro} satisfies:
\begin{equation}\label{eq:B-reduction}
  \left|
    \frac{c^2}{T^2}\frac{\int_{-T}^T x^2\psi_m^2\psi_n^2\,dx}
    {\int_{-T}^T\psi_m^2\psi_n^2\,dx}
    - \frac{c^2}{T^2}\frac{\int_{-T}^T x^2\psi_m^2\,dx}
      {\lambda_m(c)}
  \right|
  \le \frac{c^2}{T^2}\cdot\frac{C_\gamma\,\omega_n(C_\gamma/\lambda_m(c))}
      {Z_{mn}},
\end{equation}
where $Z_{mn} := \int_{-T}^T\psi_m^2\psi_n^2\,dx / \lambda_m(c)\lambda_n(c)$.
An analogous bound holds for the kinetic energy term.
\end{lemma}

\begin{proof}
The function $f_m(x) := x^2\psi_m(x)^2/\lambda_m(c)$
lies in $C^1([-T,T])$ since $\psi_m \in C^\infty$.
We bound $\|f_m'\|_{L^\infty}$ via the ORX pointwise estimates
\cite[Prop.~2.5, Prop.~2.7]{OsipovRokhlinXiao}:
\[
  \|\psi_m\|_{L^\infty} \le C(1+\chi_m)^{1/4},
  \qquad
  \|\psi_m'\|_{L^\infty} \lesssim \chi_m^{3/4}.
\]
The derivative bound uses ORX Prop.~2.7 ($\|\psi_m'\|_{L^2}^2 \sim \chi_m$
via the SL energy identity~\eqref{eq:SL-energy}), the 1D Sobolev
interpolation $\|u'\|_{L^\infty} \lesssim \|u'\|_{L^2}^{1/2}\|u''\|_{L^2}^{1/2}$,
and $\psi_m'' = O(\chi_m)\psi_m$ from the eigenvalue equation~\eqref{eq:SL}.
Since $\chi_m \le C_\gamma N^2$ and $\lambda_m(c) \ge c_\gamma > 0$
(Assumption~\ref{ass:lw}) uniformly for $m \le \gamma N$, we obtain
\[
  \|f_m'\|_{L^\infty}
  \le \frac{2T\|\psi_m\|_{L^\infty}^2
            + 2T^2\|\psi_m\|_{L^\infty}\|\psi_m'\|_{L^\infty}}
           {\lambda_m(c)}
  \le \frac{C_\gamma}{\lambda_m(c)}.
\]
Apply hypothesis~\eqref{eq:weak-conv-hyp} with $f = f_m$:
\[
  \Bigl|\int x^2\psi_m^2\psi_n^2\,dx
    - \lambda_n(c)\int x^2\psi_m^2\rhocl\,dx\Bigr|
  \le \lambda_m(c)\lambda_n(c)\,\omega_n(C_\gamma/\lambda_m(c)).
\]
The classical reference integral satisfies
$\lambda_n(c)\int x^2\psi_m^2\rhocl\,dx
\approx \lambda_n(c)\int x^2\psi_m^2\,dx / \lambda_m(c)$
by an application of the same argument with $m$ and $n$ exchanged.
Dividing by $\int\psi_m^2\psi_n^2\,dx = \lambda_m(c)\lambda_n(c)Z_{mn}$
yields~\eqref{eq:B-reduction}.

\emph{Closure.}
Hypothesis~\eqref{eq:weak-conv-hyp} is not an assumption:
it is supplied unconditionally by Paper~IV \cite{PaperIV},
Theorem~1, with modulus $\omega_n(\|f'\|_{L^\infty}) = C\|f'\|_{L^\infty}/n$
for all $f \in C^1([-T,T])$ and $n \le \gamma N$.
The present lemma therefore holds unconditionally in the bulk regime.
\end{proof}

\begin{remark}[What the reduction says]\label{rem:bulk-reduction-meaning}
Lemma~\ref{lem:bulk-reduction} reduces the energy equidistribution
question to showing that the ratio
$\int x^2\psi_m^2\psi_n^2 / \int\psi_m^2\psi_n^2$
is well-approximated by $\mathbb{E}_{\rho_m}[x^2]
= \int x^2\psi_m^2\,dx/\lambda_m(c)$.
Once $\mathbb{E}_{\rho_m}[x^2]$ is controlled via the SL energy
identity~\eqref{eq:SL-energy}, the full equidistribution
$A_{mn}+B_{mn} = \mumN/2 + o(\mumN)$ follows.
\end{remark}

\begin{remark}[Why Cauchy--Schwarz fails here]\label{rem:cs-fails}
A direct Cauchy--Schwarz bound
$|\int f\psi_m^2\psi_n^2| \le \|\psi_n\|_{L^\infty}^2\|\psi_m^2\|_{L^1}\|f\|_{L^\infty}$
gives an $O(1)$-bound on the ratio in~\eqref{eq:B-reduction},
which does not improve on the trivial $|B_{mn}| \le \chi_m(c)$.
The oscillation-averaging that makes $\rho_{mn} \approx \rho_m$
requires the weak-convergence input~\eqref{eq:weak-conv-hyp};
no pointwise $L^\infty$ bound on $\psi_n^2$ suffices.
\end{remark}

%%=================================================================
\section{Off-Diagonal Regime}\label{sec:offdiag}
%%=================================================================

\subsection{Near-eigenvector structure}\label{ssec:neareig}

\begin{lemma}[Near-eigenvector identity]\label{lem:near-eig}
For all $m,n \ge 0$:
\begin{equation}\label{eq:near-eig}
  D_c f_{mn} = \mumN f_{mn} + R_{mn},
  \qquad
  R_{mn}(x) := -p(x)\bigl[\psi_m''(x)\psi_n(x) + \psi_m(x)\psi_n''(x)\bigr]
               - p'(x)\bigl[\psi_m'(x)\psi_n(x) + \psi_m(x)\psi_n'(x)\bigr].
\end{equation}
\end{lemma}

\begin{proof}
Direct computation:
$D_c(\psi_m\psi_n)
= -(p\psi_m')'\psi_n - \psi_m(p\psi_n')' - 2p\psi_m'\psi_n'
+ (c^2x^2/T^2)\psi_m\psi_n
= \chi_m\psi_m\psi_n + \chi_n\psi_m\psi_n - 2p\psi_m'\psi_n'
= \mumN f_{mn} + R_{mn}$,
with $R_{mn} = -2p\psi_m'\psi_n'$.
Expanding:
$-2p\psi_m'\psi_n' = -(p(\psi_m\psi_n)')' + p\psi_m''\psi_n + p\psi_m\psi_n'' + p'(\psi_m'\psi_n+\psi_m\psi_n')$;
rearranging gives~\eqref{eq:near-eig}.
\end{proof}

\begin{lemma}[Spectral gap, off-diagonal]\label{lem:gap-od}
For $k \ge N$ and $|m-n|\ge\delta N$:
$|\chi_k(c) - \mumN| \ge c_3(\delta)|m-n|^2$.
\end{lemma}

\begin{proof}
From~\eqref{eq:chi-asym}: $\chi_k \ge k(k+1) \ge N^2$,
while $\mumN = \chi_m+\chi_n \le C_\delta N^2$ for $m,n \le N$.
The gap grows quadratically in $|m-n|$.
\end{proof}

\subsection{Mean spectral localization}

\begin{proposition}[Mean spectral localization]\label{prop:mean-loc}
For all $m,n \le N$ with $f_{mn}\not\equiv 0$:
\begin{equation}\label{eq:mean-loc}
  \Emn[\chi_k]
  = \mumN + E_{mn},
  \qquad
  E_{mn}
  := \frac{\ip{R_{mn}}{f_{mn}}_{L^2}}{\|f_{mn}\|_{L^2}^2},
\end{equation}
where $E_{mn} \ge 0$ (established in
Proposition~\ref{prop:spectral-lower} below).
\end{proposition}

\begin{proof}
$\ip{D_c f_{mn}}{f_{mn}} = \mumN\|f_{mn}\|^2 + \ip{R_{mn}}{f_{mn}}$.
Divide by $\|f_{mn}\|^2$.
\end{proof}

\subsection{Integration by parts and exact energy formula}

\begin{lemma}[IBP exact formula]\label{lem:ibp-exact}
\begin{equation}\label{eq:ibp-exact}
  \ip{R_{mn}}{f_{mn}}_{L^2}
  = \int_{-T}^T p(x)\bigl[(\psi_m')^2\psi_n^2 + \psi_m^2(\psi_n')^2\bigr]\,dx
  + \frac{c^2}{T^2}\int_{-T}^T x^2\psi_m^2\psi_n^2\,dx
  - \frac{\mumN}{2}\|f_{mn}\|_{L^2}^2.
\end{equation}
\end{lemma}

\begin{proof}
Integrate $\ip{R_{mn}}{f_{mn}} = -2\int p\psi_m'\psi_n'\psi_m\psi_n\,dx$
by parts twice, using $p(\pm T)=0$. The kinetic cross-term becomes
$\int p(\psi_m')^2\psi_n^2 + \int p\psi_m^2(\psi_n')^2$
minus the eigenvalue contributions, which combine to give
$-(\mumN/2)\|f_{mn}\|^2$.
\end{proof}

\begin{corollary}[Explicit energy formula]\label{cor:emn-explicit}
\begin{equation}\label{eq:emn-explicit}
  \Emn[\chi_k]
  = \frac{\mumN}{2}
  + \frac{
    \frac{1}{2}\int_{-T}^T p(x)\bigl[(\psi_m')^2\psi_n^2
      + \psi_m^2(\psi_n')^2\bigr]\,dx
    + \frac{c^2}{T^2}\int_{-T}^T x^2\psi_m^2\psi_n^2\,dx
    }{\|f_{mn}\|_{L^2}^2}.
\end{equation}
\end{corollary}

\begin{proposition}[Unconditional spectral lower bound]\label{prop:spectral-lower}
For all $m,n \le N$ with $f_{mn}\not\equiv 0$:
$\Emn[\chi_k] \ge \mumN/2$.
\end{proposition}

\begin{proof}
Every term in the numerator of~\eqref{eq:emn-explicit} is non-negative
($p \ge 0$ on $(-T,T)$, $x^2 \ge 0$).
\end{proof}

\subsection{Unconditional uniform off-diagonal bound}

\begin{theorem}[Off-diagonal uniform bound]\label{thm:offdiag}
For $|m-n| \ge \delta N$ and all $c \ge 1$:
$\|(I-\PN)f_{mn}\|_{L^2} \le C_\delta T^{1/2}$.
\end{theorem}

\begin{proof}
$\|(I-\PN)f_{mn}\|_{L^2} \le \|f_{mn}\|_{L^2}
\le \|\psi_m\|_{L^4}\|\psi_n\|_{L^4}
\le C_\delta T^{1/2}$
by H\"{o}lder and $\|\psi_n\|_{L^4} \le C T^{1/4}$
(from $\|\psi_n\|_{L^2} = \lambda_n^{1/2} \le 1$
and $\|\psi_n\|_{L^\infty} \le C(1+\chi_n)^{1/4}$
by ORX Prop.~2.5).
\end{proof}

\subsection{Conditional off-diagonal decay}

\begin{theorem}[Off-diagonal algebraic decay]\label{thm:offdiag-strong}
Under Assumption~\ref{ass:specloc}, for $|m-n| \ge \delta N$:
$\|(I-\PN)f_{mn}\|_{L^2} \le C_p(1+|m-n|)^{-p}$.
\end{theorem}

\begin{proof}
By Assumption~\ref{ass:specloc} and Lemma~\ref{lem:gap-od}:
\[
  \|(I-\PN)f_{mn}\|_{L^2}^2
  = \sum_{k\ge N}|a_k^{mn}|^2
  \le C_p^2\sum_{k\ge N}(1+c_3(\delta)|m-n|^2)^{-2p}
  \le \frac{C_p'}{(1+|m-n|)^{2p}},
\]
since the sum converges for $2p > 1$.
\end{proof}

\subsection{Product structure and spectral gap mechanism}

\begin{remark}[Product localization via spectral gap]\label{rem:product-localization}
The near-eigenvector identity~\eqref{eq:near-eig} gives the
algebraic decomposition
\[
  D_c f_{mn} = (\chi_m + \chi_n)f_{mn} + R_{mn}.
\]
In the bulk regime $m,n \le \gamma N$, the eigenvalue sum satisfies
$\chi_m + \chi_n \le 2C_\gamma N^2$,
while for $k \ge N$ we have $\chi_k \ge N^2$.
The spectral gap
\[
  \chi_N - (\chi_m + \chi_n)
  \ge N^2 - 2C_\gamma N^2
  = N^2(1 - 2C_\gamma)
  \gg 0
\]
holds for $\gamma < 1/2$ and all sufficiently large $N$,
using the precise asymptotics
$\chi_n(c) = n(n+1) + \tfrac{1}{2}c^2 + O(n^{-1}c^2)$
from~\eqref{eq:chi-asym} with $\chi_n \sim n^2$ for $n \to \infty$.
This gap --- not any microlocal WKB argument --- is the mechanism
that forces concentration of $f_{mn}$ below index $N$:
the remainder $R_{mn}$ is spectrally small relative to
$\chi_N - \mumN$ in the bulk.
\end{remark}

%%=================================================================
\section{Edge Regime}\label{sec:edge}
%%=================================================================

\begin{proposition}[Bandwidth doubling obstruction]\label{prop:bwdoubling}
For $n \sim N$ (i.e.\ $n \ge (1-\delta)N$ for any fixed $\delta > 0$):
$\Eout(f_{nn}) \ge c_0 > 0$,
where $c_0$ depends only on $\delta$ and $c$.
\end{proposition}

\begin{proof}
For $n$ near $N$, the PSWF $\psi_n$ has significant Fourier mass
near $\pm\omega$ (this is the regime of weak concentration,
$\lambda_n(c) \to 0$ as $n \to \infty$ relative to $N$).
The convolution $\widehat{f_{nn}} = |\hat\psi_n|^2 * |\hat\psi_n|^2$
therefore has non-trivial mass in $(\omega,2\omega)$:
\[
  \Eout(f_{nn})
  = \int_{|\xi|>\omega}|\hat f_{nn}(\xi)|^2\,d\xi
  \ge \int_\omega^{2\omega}|\hat f_{nn}(\xi)|^2\,d\xi
  \ge c_0(\delta,c) > 0.
\]
This is a consequence of the band-doubling support
$\operatorname{supp}\widehat{f_{nn}} \subset [-2\omega,2\omega]$
combined with the non-vanishing of $\hat\psi_n$ near $\pm\omega$
in the edge regime.
\end{proof}

\begin{remark}[Edge as different phase, not harder bulk]\label{rem:edge-structural}
The obstruction in Proposition~\ref{prop:bwdoubling} is not a
failure of the bulk method at a difficult parameter range:
it reflects a genuinely different phase of the analysis.
In the bulk, $\chi_m + \chi_n \ll \chi_N$ provides spectral
separation; in the edge regime this gap collapses:
$\chi_m + \chi_n \sim 2\chi_N \gg \chi_N$.
Consequently both the Slepian-identity approach and the
spectral-gap mechanism fail simultaneously.
The edge requires different tools (e.g.\ Airy-function
asymptotics or a refined WKB analysis near the classical turning
points at $x = \pm T$); this is left as a separate program.
\end{remark}

%%=================================================================
\section{Open Problems}\label{sec:open}
%%=================================================================

\begin{problem}[Bulk convolution decay]\label{prob:bulkconv}
Prove Assumption~\ref{ass:bulkconv}: for $m,n \le \gamma N$,
$\Eout(f_{mn}) \le C_\gamma e^{-\alpha_\gamma c}$.
\end{problem}

\begin{problem}[Weak convergence of PSWF densities]\label{prob:pswf-weak-limit}
Prove that $\psi_n^2\,dx \rightharpoonup \lambda_n(c)\,\rhocl\,dx$
with explicit rate as $n \to \infty$.
This problem is resolved unconditionally by Paper~IV
\cite{PaperIV}, Theorem~1.
\end{problem}

\begin{problem}[Refined commutator bounds]\label{prob:comm-refined}
Establish $A_{mn}+B_{mn} = \mumN/2 + o(\mumN)$
for $m,n \le \gamma N$ without any WKB assumption.
By Lemma~\ref{lem:bulk-reduction} and Paper~IV \cite{PaperIV},
Theorem~1, this reduces to controlling $Z_{mn}$ from below,
which is a second-moment problem for the product density.
\end{problem}

\begin{problem}[PSWF Clebsch--Gordan]\label{prob:cg}
Prove Assumption~\ref{ass:specloc}: establish pointwise
spectral localization of $a_k^{mn}$ around $\chi_k \approx \mumN$.
\end{problem}

\begin{problem}[Edge regime]\label{prob:edge}
Analyse the regime $m,n \ge (1-\delta)N$.
Proposition~\ref{prop:bwdoubling} shows the Slepian--gap
mechanism fails here; a separate method is required.
\end{problem}

%%=================================================================
\appendix
%%=================================================================

\section{ORX Pointwise Estimates Used}\label{app:orx}

We collect the ORX estimates \cite{OsipovRokhlinXiao} used
in the body of this paper.

\begin{itemize}
\item \textbf{Prop.~2.5} ($L^\infty$ bound):
  $\|\psi_n^{(c)}\|_{L^\infty([-T,T])} \le C(1+\chi_n(c))^{1/4}$.
\item \textbf{Prop.~2.7} ($L^2$ derivative bound):
  $\|\bigl(\psi_n^{(c)}\bigr)'\|_{L^2([-T,T])}^2 \lesssim \chi_n(c)$,
  which follows directly from the SL energy identity~\eqref{eq:SL-energy}
  and $\lambda_n(c) \le 1$.
\item \textbf{1D Sobolev interpolation}: For $u \in H^2([-T,T])$:
  $\|u'\|_{L^\infty} \lesssim \|u'\|_{L^2}^{1/2}\|u''\|_{L^2}^{1/2}$.
  Combined with $\psi_n'' = O(\chi_n)\psi_n$ from~\eqref{eq:SL},
  this gives $\|\psi_n'\|_{L^\infty} \lesssim \chi_n^{3/4}$.
\end{itemize}

These three ingredients, together with the eigenvalue
asymptotics~\eqref{eq:chi-asym}, are the only analytic inputs
to Lemma~\ref{lem:bulk-reduction}.
No WKB expansion and no result from Paper~IV is used in
deriving the $L^\infty$ bounds themselves.

%%=================================================================
\begin{thebibliography}{99}

\bibitem{PaperI}
Anonymous Author.
\newblock Frame coercivity of prolate Gram forms.
\newblock {\em Paper~I in the Prolate--Weil Program}, 2024.

\bibitem{PaperII}
Anonymous Author.
\newblock Spectral theory of the prolate concentration operator.
\newblock {\em Paper~II in the Prolate--Weil Program}, 2024.

\bibitem{PaperII_quad}
Anonymous Author.
\newblock Quadrature compatibility and the Prolate--Weil conjecture.
\newblock {\em Paper~II (quadrature supplement) in the Prolate--Weil Program},
  2024.

\bibitem{PaperIV}
Anonymous Author.
\newblock Semiclassical analysis of PSWF densities: weak convergence to
  the classical equilibrium measure.
\newblock {\em Paper~IV in the Prolate--Weil Program}, 2025.

\bibitem{OsipovRokhlinXiao}
A.~Osipov, V.~Rokhlin, and H.~Xiao.
\newblock {\em Prolate Spheroidal Wave Functions of Order Zero}.
\newblock Springer, New York, 2013.

\bibitem{Widom1964}
H.~Widom.
\newblock Asymptotic behavior of the eigenvalues of certain integral equations
  II.
\newblock {\em Arch.\ Rational Mech.\ Anal.}, 17:215--229, 1964.

\bibitem{Slepian1978}
D.~Slepian.
\newblock Prolate spheroidal wave functions, Fourier analysis and uncertainty
  V: The discrete case.
\newblock {\em Bell Syst.\ Tech.\ J.}, 57:1371--1430, 1978.

\bibitem{Weil1952}
A.~Weil.
\newblock Sur les ``formules explicites'' de la th\'{e}orie des nombres
  premiers.
\newblock {\em Comm.\ S\'{e}m.\ Math.\ Univ.\ Lund}, 252--265, 1952.

\bibitem{CCM2025}
A.~Connes, H.~Consani, and C.~Marcolli.
\newblock Weil positivity and trace formula.
\newblock {\em Preprint}, 2025.

\end{thebibliography}

\end{document}
